%% kapa: exam=1 round=31 year=2020 period=2 form=A subject=accounting subject_ko=회계학 seq=10 number=50
%% source: Q-Net 감정평가사 기출문제 (cst00309) artlSeq=5207219 file=extracted/1008/5207219/A형/A형/2020년도 제31회 감정평가사 1차 시험_2교시_A형.pdf
%% conversion: kordoc 4.24.51 -> tools/gen1.py (auto); review=proofread-sonnet; flags=none
\begin{question}{50}
(주)감평은 신약개발을 위해 20x1년 중에 연구활동관련 ￦500,000, 개발활동관련 ￦800,000을 지출하였다. 개발활동에 소요된 ￦800,000 중 ￦300,000은 20x1년 3월 1일부터 동년 9월 30일까지 지출되었으며 나머지 금액은 10월 1일 이후에 지출되었다. (주)감평의 개발활동이 무형자산 인식기준을 충족한 것은 20x1년 10월 1일부터이며, (주)감평은 20x2년 초부터 20x2년 말까지 ￦400,000을 추가 지출하고 신약개발을 완료하였다. 무형자산으로 인식한 개발비는 20x3년 1월 1일부터 사용이 가능하며, 내용연수 4년, 잔존가치 ￦0, 정액법으로 상각하고, 원가모형을 적용한다. (주)감평의 20x3년 개발비 상각액은?
\begin{choices}
  \item ￦225,000
  \item ￦250,000
  \item ￦300,000
  \item ￦325,000
  \item ￦350,000
\end{choices}
\end{question}
\answer{1}
